\section{Nghiệm bình phương tối thiểu (Analytical Least-Squares Solution)}

Bên cạnh phương pháp lặp Gradient Descent (GD) dựa trên đạo hàm, bài toán hồi quy tuyến tính còn có thể được giải quyết triệt để bằng phương pháp giải tích.

\subsection{Cơ sở toán học: Phương trình chuẩn (Normal Equations)}
Mục tiêu của tối ưu hóa là tìm điểm cực tiểu của hàm chi phí $J(\theta)$. Theo nguyên lý tích phân, tại điểm cực tiểu, đạo hàm của hàm chi phí bằng 0[span_0](start_span)[span_0](end_span):
\begin{equation}
\nabla_{\theta} J(\theta) = 0 \implies \theta_{LS} = (X^T X)^{-1} X^T y
\end{equation}

\noindent\textbf{Trong đó:}
\begin{itemize}
    \item $X$: Mã trận thiết kế chứa các đặc trưng đầu vào (kích thước $m \times (n+1)$).
    \item $y$: Vector cột chứa các giá trị thực tế của biến mục tiêu.
    \item $\theta_{LS}$: Vector tham số tối ưu toán học[span_1](start_span)[span_1](end_span).
    \item $(X^T X)^{-1}$: Ma trận nghịch đảo của tích $X^T X$.
\end{itemize}
Việc tính toán $\theta_{LS}$ thông qua phương trình này mang lại nghiệm chính xác ngay trong một bước tính toán duy nhất, không cần lặp.

\subsection{Kết quả thực nghiệm và So sánh tham số}
Để kiểm chứng tính chính xác của thuật toán Gradient Descent từ cài đặt ở phần trước, nhóm tiến hành tính toán nghiệm giải tích[span_2](start_span)[span_2](end_span):

\begin{itemize}
    \item \textbf{Về bộ tham số tối ưu ($\theta$):}
    \begin{itemize}
        \item Nghiệm từ Gradient Descent (sau $1000$ vòng lặp): $\theta_0 \approx 12.3131$, $\theta_1 \approx 0.8214$[span_3](start_span)[span_3](end_span)
        \item Nghiệm Giải tích (Nghiệm đúng): $\theta_0 \approx 12.3146$, $\theta_1 \approx 0.8309$[span_4](start_span)[span_4](end_span)
    \end{itemize}
    \item \textbf{Về sai lệch hàm chi phí (Cost Difference = $1.1\text{e}-06$):}
    \begin{itemize}
        \item Hệ thống ghi nhận sự chênh lệch giá trị Cost $J(\theta)$ giữa hai phương pháp là $1.1\text{e}-06$ (tức là $0.000001183$)[span_5](start_span)[span_5](end_span).
    \end{itemize}
    \item \textbf{Nguyên nhân và Kết luận:} Mức chênh lệch $1.1 \times 10^{-6}$ này là một giá trị cực kỳ nhỏ (gần bằng $0$). Điều này chứng tỏ:
    \begin{enumerate}
        \item Thuật toán Gradient Descent mà nhóm tự lập từ đầu (scratch) hoạt động hoàn toàn chính xác[span_6](start_span)[span_6](end_span).
        \item Các siêu tham số (Hyperparameters) được cấu hình trong GD bao gồm learning rate $\gamma = 0.01$ và giới hạn $1000$ vòng lặp đã đạt hiệu quả hội tụ tối ưu[span_7](start_span)[span_7](end_span).
    \end{enumerate}
\end{itemize}
